\begin{document}
\section*{Introducción a las probabilidades }

\subsubsection*{Experimiento aleatorio}
Es cualquier experimento cuyo resultado no se puede predecir antes de realizar el experimento porque consta con más de un resultado posible.

\subparagraph*{Ejemplos} -- 
\begin{itemize}
\item{Lanzar una moneda normal sobre una superficie plana y observar la cara superior. Puede ocurrir cara o sello.}
\item{Lanzar un dado sobre una superficie plana y observar la parte superior. Puede ocurrir que aparezca uno de los siguientes números : 1 .. 6}
\item{Extraer una bola de una urna que contiene bolas de diferentes colores.}
\end{itemize}

\subsubsection*{Espacio de muestra}
Es el conjunto formado por todos los resultados posibles de un experimento aleatorio. A su vez esté se comporta como el conjunto universal.

\subparagraph*{Ejemplos} 
Para el experimento E1, su espacio muestral es $S_1 = \{C, S\}$
Para el experimento E2, su espacio muestral  es $S_2 = \{1; 2 ; 3; 4;5 ;6 \}$


\end{document}